\chapter{Sieves and the Parity Problem}
\begin{center}\itshape
A sieve counts what survives once the multiples of small primes are struck out.\\
It can count almost-primes, bound primes from above, and do it uniformly everywhere.\\
It cannot, on its own, tell a prime from a product of two primes: that is the parity problem.
\end{center}

In 1915 a young Norwegian, Viggo Brun, published a short paper on twin primes and Goldbach's problem, and over the next five years turned Eratosthenes' sieve into a working tool. Legendre's inclusion--exclusion (Volume~5, Proposition~8.1) has $2^{\pi(z)}$ terms and collapses under its own error terms. Brun's idea was to \emph{truncate} it, keeping only the divisors with few prime factors, and to control the truncation by Bonferroni's inequalities. With it he proved in 1919 that the sum of the reciprocals of the twin primes converges (Volume~5, Corollary~8.6), and in 1920 that every large even number is a sum of two integers with at most nine prime factors each.

Selberg's $\Lambda^2$ sieve of 1947 (Volume~5, Theorem~8.2) replaced Brun's combinatorics by a quadratic form, and gave the Brun--Titchmarsh inequality with the constant $2$. Selberg also saw what no sieve can do. In his lectures on sieves he gave a pair of sequences, the integers with an even and with an odd number of prime factors, which look exactly alike to a sieve but of which one contains all the primes and the other none. So the sieve, by itself, cannot produce primes.

This chapter proves the \emph{fundamental lemma} of sieve theory, Brun's method in its sharp modern form: when the sifting primes are small compared with the level, the sieve is an asymptotic formula, from above and from below. It proves Selberg's parity example, and with it that the constant $2$ in Brun--Titchmarsh is the best a sieve can give. And it uses the sieve where it shines, as one of three tools in Hooley's proof that the generalised Riemann hypothesis implies Artin's conjecture on primitive roots.

\begin{reminder}[from Volumes 5 and 12]
\emph{The sieve setting} (Volume~5, Chapter~8). $\mathcal A$ is a finite list of positive integers, $\mathcal P$ a set of primes, $z\ge2$, $P(z)=\prod_{p\in\mathcal P,\,p\le z}p$, and $S(\mathcal A,\mathcal P,z)$ the number of $a\in\mathcal A$ coprime to $P(z)$. For $d\mid P(z)$, $\mathcal A_d$ is the sublist of multiples of $d$, and
$|\mathcal A_d|=g(d)X+r_d$ with $g$ multiplicative. Selberg's sieve (Theorem~8.2) gives upper bounds; Brun--Titchmarsh (Theorem~8.4) says $\pi(x+N)-\pi(x)\le\frac{2N}{\log N}(1+o(1))$. \emph{Mertens} (Theorems~1.5 and~1.6): $\sum_{p\le x}\frac1p=\log\log x+M+O(\frac1{\log x})$ and $\prod_{p\le x}(1-\frac1p)\sim\frac{e^{-\gamma}}{\log x}$; also $\prod_{p\le x}p\le4^x$ (Proposition~1.2).
\emph{Decomposition groups} (Volume~12, Theorem~2.18). If $L/\Q$ is Galois and $P$ a prime of $L$ above $p$, the decomposition group $D_P$ maps onto $\operatorname{Gal}(k_P/\mathbb F_p)$ with kernel the inertia group $I_P=\{\sigma\in D_P:\sigma(x)\equiv x\pmod P\text{ for all }x\in\mathcal O_L\}$, of order $e$; $|D_P|=ef$, and when $e=1$, $D_P$ is generated by the Frobenius, $\operatorname{Frob}_P(x)\equiv x^p\pmod P$. So $p$ splits completely in $L$ exactly when $D_P$ is trivial.
\end{reminder}

\section{The fundamental lemma}

\subsection*{Brun's pure sieve}

Write $\omega(d)$ for the number of prime factors of $d$. Legendre's identity $\mathbf 1_{n=1}=\sum_{d\mid n}\mu(d)$ has as many terms as $n$ has divisors. Brun cut the sum at divisors with at most $r$ prime factors, and observed that the cut sums are alternately too large and too small.

\begin{lemma}[Bonferroni]\label{lem:bonf}
Let $n$ be squarefree, $r\ge0$, and $E_r(n)=\sum_{d\mid n,\ \omega(d)\le r}\mu(d)-\mathbf 1_{n=1}$. Then $(-1)^rE_r(n)\ge0$ and $|E_r(n)|\le\#\{d\mid n:\omega(d)=r+1\}$.
\end{lemma}
\begin{proof}
For $n=1$, $E_r(n)=0$. Otherwise let $m=\omega(n)\ge1$. Then $E_r(n)=\sum_{j=0}^r(-1)^j\binom mj=(-1)^r\binom{m-1}r$, by induction on $r$ using $\binom mj=\binom{m-1}j+\binom{m-1}{j-1}$. This has the sign $(-1)^r$. If $m\le r$ it vanishes, and if $m\ge r+1$, $\binom{m-1}r=\frac{r+1}m\binom m{r+1}\le\binom m{r+1}$, the number of divisors of $n$ with $r+1$ prime factors.
\end{proof}

The same inequality, averaged with multiplicative weights, controls the main terms.

\begin{lemma}\label{lem:bonfavg}
Let $P'$ be squarefree, $g$ multiplicative with $0\le g(p)\le1$ for $p\mid P'$, $V'=\prod_{p\mid P'}(1-g(p))$ and $G'=\sum_{p\mid P'}g(p)$. For $r\ge0$,
\[
\Bigl|\sum_{\substack{d\mid P'\\\omega(d)\le r}}\mu(d)g(d)-V'\Bigr|\le\sum_{\substack{d\mid P'\\\omega(d)=r+1}}g(d)\le\frac{G'^{\,r+1}}{(r+1)!}.
\]
\end{lemma}
\begin{proof}
For $n\mid P'$ let $w(n)=\prod_{p\mid n}g(p)\prod_{p\mid P'/n}(1-g(p))\ge0$: the probability that the random set of primes, each $p\mid P'$ chosen independently with probability $g(p)$, has product $n$. For $d\mid P'$, $\sum_{d\mid n\mid P'}w(n)=g(d)$, since the primes outside $d$ contribute $\prod(g(p)+1-g(p))=1$. Hence
\[
\sum_{\substack{d\mid P'\\\omega(d)\le r}}\mu(d)g(d)-V'=\sum_{n\mid P'}w(n)E_r(n),
\]
as $V'=w(1)$. By Lemma~\ref{lem:bonf} its absolute value is at most $\sum_nw(n)\#\{d\mid n:\omega(d)=r+1\}=\sum_{d\mid P',\,\omega(d)=r+1}g(d)$. Finally, expanding $G'^m$ produces each $g(d)$ with $d\mid P'$, $\omega(d)=m$, exactly $m!$ times, and other non-negative terms.
\end{proof}

If we sift by all primes up to $z$ at once, $G'=\sum_{p\le z}g(p)$ is about $\kappa\log\log z$ when $g(p)\approx\frac\kappa p$, so the error $\frac{G'^{r+1}}{(r+1)!}$ is smaller than the main term $V'\approx(\log z)^{-\kappa}$ only when $r$ is a large multiple of $\log\log z$. The divisors then go up to $z^r=z^{C\log\log z}$. This is Brun's \emph{pure sieve}: it works, but only for $z$ as small as $x^{1/(C\log\log x)}$. The cure, also Brun's, is to cut the primes into ranges and truncate each range separately, with more terms where the primes are small.

\begin{lemma}\label{lem:prodineq}
Let $a_j,u_j,\ell_j$ ($1\le j\le J$) be real numbers with $0\le a_j\le u_j$ and $\ell_j\le a_j$. Then
\[
\prod_ju_j\ \ge\ \prod_ja_j\ \ge\ \prod_ju_j-\sum_i(u_i-\ell_i)\prod_{j\ne i}u_j.
\]
\end{lemma}
\begin{proof}
Telescoping, $\prod_ju_j-\prod_ja_j=\sum_i(u_i-a_i)\prod_{j<i}a_j\prod_{j>i}u_j$. Each term is at least $0$, and at most $(u_i-\ell_i)\prod_{j\ne i}u_j$, because $0\le a_j\le u_j$ and $0\le u_i-a_i\le u_i-\ell_i$.
\end{proof}

\subsection*{The lemma}

Let $V(z)=\prod_{p\mid P(z)}(1-g(p))$, the proportion that would survive if the events $p\mid a$ were independent with probabilities $g(p)$. We assume, for constants $\kappa>0$ and $K\ge2$,
\[
0\le g(p)\le\min\Bigl(\frac\kappa p,\,1-\frac1K\Bigr)\qquad(p\mid P(z)).\tag{$\Omega$}
\]
The number $\kappa$ is the \emph{dimension} of the sieve: about $\kappa$ residue classes are removed modulo each prime. The second condition says that no prime removes everything.

\begin{theorem}[the fundamental lemma]\label{thm:fl}
Assume $(\Omega)$. There is a constant $C$, depending only on $\kappa$ and $K$, such that for $z\ge2$, $s\ge1$ and $D=z^s$,
\[
S(\mathcal A,\mathcal P,z)=XV(z)\bigl(1+\theta Ce^{-s}\bigr)+\theta'\sum_{\substack{d\mid P(z)\\d<D}}\tau(d)|r_d|,
\]
with $|\theta|\le1$, $|\theta'|\le1$.
\end{theorem}

So the sieve gives the expected answer, with a relative error decaying exponentially in $s=\frac{\log D}{\log z}$, as long as the remainders $r_d$ are small on average up to the level $D$.

\begin{proof}
\emph{Constants.} By Mertens (Volume~5, Theorem~1.5) there is an absolute $B\ge1$ with $\sum_{w<p\le w'}\frac1p\le\log\frac{\log w'}{\log w}+B$ for $2\le w\le w'$, and $\sum_{p\le4}\frac1p\le B$. For $0\le g\le1-\frac1K$ we have $\log\frac1{1-g}\le2g\log K$, because $\frac1g\log\frac1{1-g}$ increases with $g$ and equals $\frac{\log K}{1-1/K}\le2\log K$ at $g=1-\frac1K$. Let $G=\kappa B$ and $v=K^{-2G}$.

\emph{Ranges.} Let $z_j=z^{2^{-j}}$ and $P_j=\prod p$ over the $p\mid P(z)$ with $z_{j+1}<p\le z_j$, for $j\ge0$. Only finitely many $P_j$ differ from $1$, and $P(z)=\prod_jP_j$. Then
\[
G_j=\sum_{p\mid P_j}g(p)\le\kappa\sum_{z_{j+1}<p\le z_j}\frac1p\le G,\qquad V_j=\prod_{p\mid P_j}(1-g(p))\ge e^{-2G_j\log K}\ge v.
\]
The sum is over $z_{j+1}<p\le z_{j+1}^2$, which is at most $\log2+B\le2B$ when $z_{j+1}\ge2$, and at most $B$ otherwise; we absorb the factor $2$ into $B$.

\emph{The two inequalities.} Let $r_j$ be even integers, chosen below. For a positive integer $n$ let $a_j(n)=\mathbf 1_{(n,P_j)=1}$ and
\[
u_j(n)=\sum_{\substack{d\mid(n,P_j)\\\omega(d)\le r_j}}\mu(d),\qquad\ell_j(n)=\sum_{\substack{d\mid(n,P_j)\\\omega(d)\le r_j-1}}\mu(d).
\]
By Lemma~\ref{lem:bonf}, applied to $(n,P_j)$ with $r_j$ even and with $r_j-1$ odd, $u_j\ge a_j\ge\ell_j$, and $u_j-\ell_j=\#\{d\mid(n,P_j):\omega(d)=r_j\}\ge0$. Since $\mathbf 1_{(n,P(z))=1}=\prod_ja_j(n)$, Lemma~\ref{lem:prodineq} gives
\[
\prod_ju_j(n)\ \ge\ \mathbf 1_{(n,P(z))=1}\ \ge\ \prod_ju_j(n)-\sum_i(u_i-\ell_i)(n)\prod_{j\ne i}u_j(n).
\]
For $d\mid P(z)$ write $d_j=(d,P_j)$, and let $\mathcal D$ be the set of $d\mid P(z)$ with $\omega(d_j)\le r_j$ for all $j$. Expanding, $\prod_ju_j(n)=\sum_{d\mid(n,P(z))}\lambda^+_d$ with $\lambda^+_d=\mu(d)\mathbf 1_{\mathcal D}(d)$. In the same way, since $\mu(d_i)=1$ when $\omega(d_i)=r_i$ is even, the right side is $\sum_{d\mid(n,P(z))}\lambda^-_d$ with
\[
\lambda^-_d=\mu(d)\mathbf 1_{\mathcal D}(d)\bigl(1-\#\{i:\omega(d_i)=r_i\}\bigr).
\]
Both are supported on $\mathcal D$, and $|\lambda^\pm_d|\le1+\omega(d)\le\tau(d)$. Summing over $a\in\mathcal A$ and using $|\mathcal A_d|=g(d)X+r_d$,
\begin{equation}\label{eq:flsandwich}
\begin{aligned}S(\mathcal A,\mathcal P,z)&\le X\sum_d\lambda^+_dg(d)+\sum_{d\in\mathcal D}\tau(d)|r_d|,\\S(\mathcal A,\mathcal P,z)&\ge X\sum_d\lambda^-_dg(d)-\sum_{d\in\mathcal D}\tau(d)|r_d|.\end{aligned}
\end{equation}

\emph{The level.} Every $d\in\mathcal D$ satisfies $d=\prod_jd_j<\prod_jz_j^{r_j}=z^{\sum_jr_j2^{-j}}$. Take $r_j=2\lceil\frac s8\rceil+2j$. Then $r_j\le\frac s4+2+2j$, and
\[
\sum_{j\ge0}r_j2^{-j}\le2\Bigl(\frac s4+2\Bigr)+2\sum_jj2^{-j}=\frac s2+8\le s\qquad\text{if }s\ge16.
\]
So for $s\ge16$, $\mathcal D\subseteq[1,D[$.

\emph{Main terms.} By multiplicativity,
\[
\begin{aligned}
\sum_d\lambda^+_dg(d)&=\prod_jU_j,\\
U_j&=\sum_{\substack{d\mid P_j\\\omega(d)\le r_j}}\mu(d)g(d).
\end{aligned}
\]
Also,
\[
\sum_d\lambda^-_dg(d)=\prod_jU_j-\sum_i\Delta_i\prod_{j\ne i}U_j,\qquad\Delta_i=\sum_{\substack{d\mid P_i\\\omega(d)=r_i}}g(d).
\]
Lemma~\ref{lem:bonfavg} gives $U_j=V_j(1+\varepsilon_j)$ with $|\varepsilon_j|\le\frac1v\frac{G^{r_j+1}}{(r_j+1)!}$, and $\Delta_i=V_i\delta_i$ with $0\le\delta_i\le\frac1v\frac{G^{r_i}}{r_i!}$. Since $m!\ge(\frac me)^m$, $\frac{G^m}{m!}\le(\frac{eG}m)^m\le e^{-4m}$ whenever $m\ge e^5G$. All the $m$ here are at least $r_j\ge\frac s4+2j$; so if $s\ge4e^5G$,
\[
\sum_j\bigl(|\varepsilon_j|+\delta_j\bigr)\le\frac2v\sum_{j\ge0}e^{-s-8j}\le\frac3ve^{-s}=:\varepsilon.
\]
Suppose also $\varepsilon\le\frac12$. Then $|\prod_j(1+\varepsilon_j)-1|\le e^{\sum|\varepsilon_j|}-1\le2\varepsilon$, so $\prod_jU_j=V(z)(1+2\theta\varepsilon)$, as $V(z)=\prod_jV_j$. And $\sum_i\Delta_i\prod_{j\ne i}U_j\le V(z)\sum_i\delta_i\prod_{j\ne i}(1+|\varepsilon_j|)\le e^{1/2}V(z)\varepsilon\le2\varepsilon V(z)$. So both main terms in \eqref{eq:flsandwich} are $V(z)(1+4\theta\varepsilon)$. This proves the theorem for $s\ge s_0=\max(16,4e^5G,\log\frac6v)$, with $C=\frac{12}v$.

\emph{Small $s$.} Let $1\le s<s_0$. The lower bound is trivial once $C\ge e^{s_0}$, since $S\ge0$. For the upper bound let $z'=D^{1/s_0}<z$. Removing fewer primes can only increase the count, so $S(\mathcal A,\mathcal P,z)\le S(\mathcal A,\mathcal P,z')$. If $z'\ge2$, the case already proved, at $(z',s_0)$, gives $S\le2XV(z')+\sum_{d\mid P(z'),\,d<D}\tau(d)|r_d|$, and
\[
\frac{V(z')}{V(z)}=\prod_{\substack{p\mid P(z)\\z'<p\le z}}\frac1{1-g(p)}\le\exp\Bigl(2\kappa\log K\sum_{z'<p\le z}\frac1p\Bigr)\le\exp\bigl(2\kappa\log K(\log s_0+B)\bigr),
\]
as $\frac{\log z}{\log z'}=\frac{s_0}s\le s_0$. If $z'<2$, then $S\le|\mathcal A|=X+r_1$, and $z<2^{s_0}$, so $V(z)\ge\exp(-2\kappa\log K\sum_{p<2^{s_0}}\frac1p)$ is bounded below. In both cases $S\le XV(z)(1+C'e^{-s})+\sum_{d\mid P(z),\,d<D}\tau(d)|r_d|$ with $C'$ depending only on $\kappa$ and $K$.
\end{proof}

\begin{think}
The weights $\lambda^\pm_d$ do not depend on $\mathcal A$ at all, only on $P(z)$ and $s$. Why does that matter for the applications, where the same weights are used for many sequences at once?
\end{think}

Because the weights are fixed in advance, the remainder is a fixed linear combination of the $r_d$, and only an \emph{average} of the $|r_d|$ over $d<D$ is needed, the quantity that Chapter~4's Bombieri--Vinogradov theorem controls for the primes.

\subsection*{First consequences}

\begin{corollary}[rough numbers]\label{cor:rough}
For $x\ge4$ and $2\le z\le\sqrt x$, let $\Phi(x,z)$ be the number of $n\le x$ with no prime factor $\le z$, and $u=\frac{\log x}{\log z}$. Then
\[
\Phi(x,z)=x\prod_{p\le z}\Bigl(1-\frac1p\Bigr)\bigl(1+O(e^{-u/2})\bigr)+O(\sqrt x\log x).
\]
In particular $\Phi(x,z)\sim e^{-\gamma}\frac x{\log z}$ when $z\to\infty$ and $u\to\infty$.
\end{corollary}
\begin{proof}
Take $\mathcal A=\{1,\dots,\lfloor x\rfloor\}$, $\mathcal P$ all primes, $X=x$, $g(d)=\frac1d$ and $|r_d|\le1$; $(\Omega)$ holds with $\kappa=1$, $K=2$. With $D=\sqrt x$, $s=\frac u2\ge1$, and $\sum_{d<D}\tau(d)=\sum_{a<D}\#\{b:ab<D\}\le D(1+\log D)$. The last statement is Mertens' Theorem~1.6.
\end{proof}

\begin{figure}[ht]
\centering
\fig{ch03-rough}
\caption{The count $\Phi(x,z)$ of $n\le10^7$ free of prime factors up to $z$, divided by the sieve prediction $x\prod_{p\le z}(1-\frac1p)$, against $u=\frac{\log x}{\log z}$: one point for each prime $z\le\sqrt x$. The ratio is $0.9545$ at $u=2$, dips to $0.935$ near $u=2.1$, overshoots to $1.005$ near $u=3$, and is within $3\cdot10^{-4}$ of $1$ from $u=3.5$ on, as the fundamental lemma predicts. The dashed curve is the limit as $x\to\infty$ on $2\le u\le3$, $e^\gamma\omega(u)=e^\gamma\frac{1+\log(u-1)}u$, from Buchstab's function $\omega$: at $u=2$ the sieve prediction is wrong by a factor $\frac{e^\gamma}2\approx0.89$, and no improvement of the error terms can fix that.}
\label{fig:rough}
\end{figure}

For $u\le2$ the prediction fails for a structural reason, visible in Figure~\ref{fig:rough}: $\Phi(x,\sqrt x)=1+\pi(x)-\pi(\sqrt x)\sim\frac x{\log x}$, while $x\prod_{p\le\sqrt x}(1-\frac1p)\sim2e^{-\gamma}\frac x{\log x}$. The independence heuristic is off by the factor $\frac{e^\gamma}2$ at $z=\sqrt x$, and the fundamental lemma needs $s$ large precisely to stay away from this.

The second consequence is Brun's theorem on almost-prime twins, in a form with an unspecified but absolute number of prime factors. Write $\Omega(n)$ for the number of prime factors of $n$ counted with multiplicity.

\begin{corollary}[Brun]\label{cor:brunalmost}
There are an integer $r$ and a constant $c>0$ such that for $x$ large,
\[
\#\{n\le x:\Omega(n)\le r\text{ and }\Omega(n+2)\le r\}\ge c\frac x{\log^2x}.
\]
\end{corollary}
\begin{proof}
Let $\mathcal A=\{n(n+2):n\le x\}$ and $\mathcal P$ all primes. As in Volume~5, Theorem~8.5, $|\mathcal A_d|=\frac{\rho(d)}dx+r_d$ for squarefree $d$, with $\rho$ multiplicative, $\rho(2)=1$, $\rho(p)=2$ for $p$ odd, and $|r_d|\le\rho(d)\le\tau(d)$. So $g(2)=\frac12$, $g(3)=\frac23$, $g(p)=\frac2p$, and $(\Omega)$ holds with $\kappa=2$, $K=3$. Let $C$ be the constant of Theorem~\ref{thm:fl} for these values and fix $s_1\ge1$ with $Ce^{-s_1}\le\frac12$. Take $D=\sqrt x$ and $z=D^{1/s_1}$.

The remainder is at most $\sum_{d<D}\tau(d)^2$. Since $\tau(n)^2\le\tau_4(n)$, the number of ways to write $n$ as a product of four factors (on prime powers, $(a+1)^2\le\binom{a+3}3$), we get $\sum_{d<D}\tau(d)^2\le D\sum_{d<D}\frac{\tau_4(d)}d\le D(1+\log D)^4$. For the main term, $1-\frac2p=(1-\frac1p)^2(1-\frac1{(p-1)^2})$, so
\[
V(z)=\frac12\prod_{3\le p\le z}\Bigl(1-\frac2p\Bigr)\ge2C_2\prod_{p\le z}\Bigl(1-\frac1p\Bigr)^2,\qquad C_2=\prod_{p>2}\Bigl(1-\frac1{(p-1)^2}\Bigr),
\]
which is $\gg\frac1{\log^2z}\asymp\frac{s_1^2}{\log^2x}$ by Mertens. Hence $S(\mathcal A,\mathcal P,z)\ge\frac12xV(z)-x^{1/2}(1+\log x)^4\ge c\frac x{\log^2x}$ for $x$ large.
Every counted $n>1$ has all the prime factors of $n$ and of $n+2$ larger than $z$. So $z^{\Omega(n)}\le n\le x$ and $z^{\Omega(n+2)}\le x+2$, and $r=\lfloor2s_1\rfloor+1$ works.
\end{proof}

The same argument handles any finite set of linear forms, and with $r$ explicit it is the start of a long race: Brun's $9$ became $2$ only with Chen (below). For upper bounds the fundamental lemma is weaker than Selberg's sieve at $z=\sqrt D$, where it is no longer asymptotic; its strength is the lower bound, and the uniformity in progressions.

\begin{corollary}[Brun--Titchmarsh in progressions]\label{cor:btap}
There is an absolute constant $C_7$ such that for $x\ge3$, $1\le q\le x^{3/4}$ and $(a,q)=1$,
\[
\pi(x;q,a)\le C_7\frac x{\varphi(q)\log x}.
\]
\end{corollary}
\begin{proof}
Let $\mathcal A$ be the $n\le x$ with $n\equiv a\pmod q$, and $\mathcal P$ the primes not dividing $q$. For $d\mid P(z)$, $(d,q)=1$, and the conditions $n\equiv a\pmod q$, $d\mid n$ define one class modulo $qd$; so $|\mathcal A_d|=\frac x{qd}+r_d$ with $|r_d|\le1$, $X=\frac xq$, $g(d)=\frac1d$. Take $D=z=(\frac xq)^{1/2}\ge x^{1/8}$, $s=1$. Then
\[
V(z)=\prod_{\substack{p\le z\\p\nmid q}}\Bigl(1-\frac1p\Bigr)\le\frac q{\varphi(q)}\prod_{p\le z}\Bigl(1-\frac1p\Bigr)\ll\frac q{\varphi(q)\log z}\ll\frac q{\varphi(q)\log x},
\]
and Theorem~\ref{thm:fl} gives $S(\mathcal A,\mathcal P,z)\ll\frac x{\varphi(q)\log x}+\sqrt{\tfrac xq}\log x$. A prime $p\equiv a\pmod q$ is either at most $z$ or counted by $S$. Finally $z+\sqrt{x/q}\log x\le2\sqrt{x/q}\log x\le\frac x{q\log x}$ for $x$ large, since $\frac{\sqrt{x/q}}{\log^2x}\ge\frac{x^{1/8}}{\log^2x}$.
\end{proof}

Volume~5, Exercise~8.4 gives the constant $2+o(1)$ with Selberg's sieve; any constant will do below.

\section{The parity problem}

\subsection*{Möbius and Liouville}

Selberg's example needs the Liouville function $\lambda(n)=(-1)^{\Omega(n)}$ to cancel with a saving of more than a power of $\log x$. The prime number theorem of Volume~5 gives $\sum_{n\le x}\mu(n)=o(x)$ (Exercise~5.3), and the classical zero-free region gives more. Since Volume~5 only proved this for $\psi$, we prove a version for $\frac1\zeta$, with a crude but sufficient exponent.

\begin{lemma}\label{lem:zetaline}
There is an absolute constant $A_5$ such that for $|t|\ge3$:
\begin{enumerate}[label=(\alph*)]
\item $|\zeta(\sigma+it)|\le A_5\log|t|$ for $\sigma\ge1-\frac1{\log|t|}$;
\item $|\zeta'(\sigma+it)|\le8A_5\log^2|t|$ for $\sigma\ge1-\frac1{4\log|t|}$.
\end{enumerate}
\end{lemma}
\begin{proof}
(a) For $\sigma\ge2$, $|\zeta(s)|\le\zeta(2)$. For $1-\frac1{\log|t|}\le\sigma\le2$ (note $\sigma>0$), apply Lemma~\ref{lem:em} with $k=2$ and $N=\lceil(|s|+2)^2\rceil\le(|t|+5)^2\le|t|^4$. For $n\le N$, $n^{-\sigma}\le\frac1nn^{1/\log|t|}\le\frac{e^4}n$, so the sum is at most $e^4(1+4\log|t|)$; also $|\frac{N^{1-s}}{s-1}|\le\frac{e^4}{|t|}$ and $|E|\le\frac32$.
(b) Apply Cauchy's estimate on the circle of radius $\frac1{4\log|t|}$ about $s$. Its points $s'$ have real part at least $1-\frac1{2\log|t|}\ge1-\frac1{\log|s'|}$, since $3\le|t|-\frac14\le|\operatorname{Im}s'|\le|t|+1\le|t|^2$, so (a) bounds $\zeta$ there by $2A_5\log|t|$.
\end{proof}

\begin{lemma}\label{lem:invzeta}
There are absolute constants $\eta\in\left]0,\frac14\right]$, $\delta_0>0$ and $A_6$ such that $\zeta$ has no zero in the region
\[
\mathcal R=\Bigl\{\sigma+it:\ \sigma\ge1-\frac\eta{\log^9|t|}\text{ if }|t|\ge3,\ \ \sigma\ge1-\delta_0\text{ if }|t|\le3\Bigr\},
\]
and $|\frac1{\zeta(s)}|\le A_6\log^9(|t|+3)$ there (with $\frac1\zeta(1)=0$).
\end{lemma}
\begin{proof}
Let $|t|\ge3$, $\ell=\log|t|$, $\delta=\frac\eta{\ell^9}\le\frac1{4\ell}$. For $\sigma\ge1+\delta$, $|\frac1{\zeta(s)}|=|\sum\mu(n)n^{-s}|\le\zeta(\sigma)\le\frac\sigma{\sigma-1}\le\frac3\delta$ (or $\zeta(2)$ if $\sigma\ge2$), which is at most $\frac3\eta\ell^9$.
At $\sigma_0=1+\delta$, the inequality $\zeta(\sigma_0)^3|\zeta(\sigma_0+it)|^4|\zeta(\sigma_0+2it)|\ge1$ of Volume~5, Theorem~5.2, with $\zeta(\sigma_0)\le\frac2\delta$ and $|\zeta(\sigma_0+2it)|\le2A_5\ell$ (Lemma~\ref{lem:zetaline}(a)), gives $|\zeta(\sigma_0+it)|\ge c_4\delta^{3/4}\ell^{-1/4}$ with $c_4=2^{-3/4}(2A_5)^{-1/4}$. For $1-\delta\le\sigma\le1+\delta$, Lemma~\ref{lem:zetaline}(b) gives
\[
|\zeta(\sigma+it)|\ge c_4\delta^{3/4}\ell^{-1/4}-16A_5\delta\ell^2=\delta^{3/4}\ell^{-1/4}\bigl(c_4-16A_5\eta^{1/4}\bigr),
\]
because $\delta^{1/4}\ell^{9/4}=\eta^{1/4}$. Choose $\eta$ with $16A_5\eta^{1/4}\le\frac{c_4}2$. Then $|\zeta(\sigma+it)|\ge\frac{c_4}2\eta^{3/4}\ell^{-7}$, so $\zeta\ne0$ and $|\frac1\zeta|\le\frac2{c_4}\eta^{-3/4}\ell^7$.
For $|t|\le3$: $\zeta$ has no zero with $\sigma\ge1$ (Volume~5, Theorem~5.2, and $\zeta(\sigma)>0$ for real $\sigma>1$), and $|\zeta|\to\infty$ at $1$. If there were zeros $\rho_n$ with $|\gamma_n|\le3$ and $\beta_n\to1^-$, a subsequence would converge to a point of $\{\sigma=1,|t|\le3\}$, which would be a zero, or the pole, by continuity; both are impossible. So some $\delta_0>0$ works, and $\frac1\zeta$, continuous on the compact set $\{1-\delta_0\le\sigma\le2,|t|\le3\}$ once we set $\frac1\zeta(1)=0$, is bounded there.
\end{proof}

\begin{proposition}\label{prop:liouville}
There is an absolute $c_5>0$ such that, for $y\ge2$,
\[
M(y)=\sum_{n\le y}\mu(n)\ll ye^{-c_5(\log y)^{1/10}},\qquad L(y)=\sum_{n\le y}\lambda(n)\ll ye^{-c_5(\log y)^{1/10}}.
\]
\end{proposition}
\begin{proof}
\emph{Perron.} For $a=1+\frac1{\log y}$, $T\ge2$ and $y$ a half-integer,
\[
M(y)=\frac1{2\pi i}\int_{a-iT}^{a+iT}\frac1{\zeta(s)}\frac{y^s}s\dd s+O\Bigl(\frac{y\log y}T+1\Bigr).
\]
Indeed, the proof of Volume~5, Proposition~7.2 integrates term by term and bounds the error of the $n$-th term by $|a_n|(\frac yn)^a\min(1,\frac1{T|\log(y/n)|})$; with $|a_n|=|\mu(n)|\le1$, the $n\le\frac y2$ and $n\ge2y$ contribute $O(\frac{y\log y}T)$, and the $n=y+h$ in between, where $|\log\frac yn|\gg\frac{|h|}y$, contribute $\sum_h\min(1,\frac y{T|h|})=O(1+\frac{y\log y}T)$. Half-integers suffice: $M$ changes by at most $1$ between them.

\emph{Moving the line.} Let $\log T=(\log y)^{1/10}$ and $\sigma_1=1-\frac\eta{\log^9T}$; for $y$ large, $\sigma_1\ge1-\delta_0$, so the rectangle with vertices $a\pm iT$, $\sigma_1\pm iT$ lies in $\mathcal R$, where $\frac1\zeta$ is holomorphic (the pole of $\zeta$ is a zero of $\frac1\zeta$) and at most $A_6\log^9(T+3)$. The integrand has no pole inside ($s=0$ is outside), so the integral equals the integral over the other three sides, which is
\[
\ll\log^9T\Bigl(y^{\sigma_1}\int_{-T}^T\frac{\mathrm dt}{|\sigma_1+it|}+\frac{2ey}T\Bigr)\ll\log^{10}T\bigl(y^{\sigma_1}+\tfrac yT\bigr).
\]
Now $\log^{10}T=\log y$, and
\[
y^{\sigma_1}=y\exp\Bigl(-\eta\frac{\log y}{\log^9T}\Bigr)=ye^{-\eta(\log y)^{1/10}},\qquad\frac yT=ye^{-(\log y)^{1/10}}.
\] So $M(y)\ll y\log y\,e^{-\eta(\log y)^{1/10}}$, which gives the bound for $M$ with any $c_5<\eta$.

\emph{Liouville.} $\sum\lambda(n)n^{-s}=\frac{\zeta(2s)}{\zeta(s)}$, so $\lambda(n)=\sum_{m^2k=n}\mu(k)$ and $L(y)=\sum_{m\le\sqrt y}M(\frac y{m^2})$. The $m\le y^{1/4}$ have $\frac y{m^2}\ge\sqrt y$ and contribute $\ll\sum_m\frac y{m^2}e^{-c_5(\frac12\log y)^{1/10}}$; the $m>y^{1/4}$ contribute at most $\sum_{m>y^{1/4}}\frac y{m^2}\le2y^{3/4}$. Adjusting $c_5$ gives the claim.
\end{proof}

The exponent $\frac1{10}$ is far from the truth (the classical region gives $\frac12$, as in Volume~5 for $\psi$), but any saving of more than a power of $\log y$ is all we need.

\subsection*{Selberg's example}

\begin{theorem}[Selberg's parity example]\label{thm:parity}
For $x\ge3$ let $\mathcal A^+$ and $\mathcal A^-$ be the lists of $n\le x$ with $\lambda(n)=1$ and $\lambda(n)=-1$. Then for every $d\ge1$,
\[
|\mathcal A^\pm_d|=\frac x{2d}+r^\pm_d,\qquad|r^\pm_d|\le1+\tfrac12\bigl|L(\tfrac xd)\bigr|,
\]
and for every $\varepsilon\in\left]0,1\right[$ and $D\le x^{1-\varepsilon}$,
\[
\sum_{d<D}\tau(d)\bigl(|r^+_d|+|r^-_d|\bigr)\ll_\varepsilon x\exp\bigl(-c_6(\varepsilon)(\log x)^{1/10}\bigr).
\]
Yet, with $\mathcal P$ all primes,
\[
S(\mathcal A^+,\mathcal P,\sqrt x)=1,\qquad S(\mathcal A^-,\mathcal P,\sqrt x)=\pi(x)-\pi(\sqrt x).
\]
\end{theorem}
\begin{proof}
Since $\lambda$ is completely multiplicative, $|\mathcal A^\pm_d|=\sum_{m\le x/d}\frac{1\pm\lambda(dm)}2=\frac12\lfloor\frac xd\rfloor\pm\frac12\lambda(d)L(\frac xd)$. For $d<D\le x^{1-\varepsilon}$, $\frac xd\ge x^\varepsilon$, and Proposition~\ref{prop:liouville} gives
\[
\sum_{d<D}\tau(d)\bigl|L(\tfrac xd)\bigr|\ll x\,e^{-c_5(\varepsilon\log x)^{1/10}}\sum_{d<D}\frac{\tau(d)}d\le x(1+\log x)^2e^{-c_5\varepsilon^{1/10}(\log x)^{1/10}},
\]
while $\sum_{d<D}2\tau(d)\le2D(1+\log D)\ll x^{1-\varepsilon}\log x$. For the last claim: an $n\le x$ with no prime factor at most $\sqrt x$ is $1$ or a prime, since a product of two primes larger than $\sqrt x$ exceeds $x$; and $\lambda(1)=1$, $\lambda(p)=-1$.
\end{proof}

So $\mathcal A^+$ and $\mathcal A^-$ have the same density $X=\frac x2$, the same local densities $g(d)=\frac1d$, and remainders that are negligible on average up to $x^{1-\varepsilon}$, almost the whole range. Any theorem that uses only this information must give the same answer for both. But one contains no primes beyond $\sqrt x$ and the other contains all of them.

\begin{corollary}[what a sieve cannot do]\label{cor:parity}
Fix $\varepsilon>0$, and suppose that for some functions $f(x)\le F(x)$, every list $\mathcal A$ of integers in $[1,x]$ with $|\mathcal A_d|=\frac Xd+r_d$ satisfies
\[
Xf(x)-R(\mathcal A)\le S(\mathcal A,\mathcal P,\sqrt x)\le XF(x)+R(\mathcal A),\qquad R(\mathcal A)=\sum_{d<x^{1-\varepsilon}}\tau(d)|r_d|.
\]
Then $f(x)\le o\bigl(\frac1{\log x}\bigr)$ and $F(x)\ge\frac{2+o(1)}{\log x}$.
\end{corollary}
\begin{proof}
Apply the lower bound to $\mathcal A^+$ and the upper bound to $\mathcal A^-$, with $X=\frac x2$. By Theorem~\ref{thm:parity}, $R(\mathcal A^\pm)=o(\frac x{\log^2x})$; so $\frac x2f(x)\le1+o(\frac x{\log x})$ and $\frac x2F(x)\ge\pi(x)-\pi(\sqrt x)-o(\frac x{\log x})$, and the prime number theorem finishes.
\end{proof}

A sieve with level of distribution below $x$ cannot prove that any list contains a prime, and its upper bound for the primes up to $x$ must be at least $2\frac X{\log x}$. For the full list $\{1,\dots,\lfloor x\rfloor\}$, write it as $\mathcal A^+$ followed by $\mathcal A^-$, and replace it by $\mathcal A^-$ listed twice: the data are unchanged up to admissible remainders, and the count of primes doubles to $\frac{2x}{\log x}$. This is the Brun--Titchmarsh bound $\frac{2N}{\log N}$ with $N=x$: the constant $2$ is optimal for sieve methods.

\begin{think}
The same example kills any sieve proof of the twin prime conjecture. Which pair of lists, in place of $\mathcal A^\pm$, would you use for twins, and which conjecture about $\lambda$ would you need to make the remainders small?
\end{think}

For twins one would use the $n$ with $\lambda(n(n+2))=\pm1$, whose remainders involve $\sum_{n\le y}\lambda(n)\lambda(n+2)$ along progressions. That this is $o(y)$ is Chowla's conjecture, still open, though Tao proved a logarithmically averaged version in 2016. So the parity obstruction for twins is believed, not proved; for the primes themselves it is a theorem.

\subsection*{Living with parity}

Three ways around the obstruction have been found. The first accepts it: a sieve can find integers with \emph{at most two} prime factors, since $\mathcal A^+$ contains many products of two primes. Chen Jingrun did this with great skill.

\begin{theorem}[Chen, 1966 and 1973; quoted]\label{thm:chen}
Every sufficiently large even number is the sum of a prime and an integer with at most two prime factors. There are infinitely many primes $p$ such that $p+2$ has at most two prime factors.
\end{theorem}

Chen sifts the list $\{N-p:p\le N\}$, so the remainders are errors for primes in progressions, and their average is controlled by the Bombieri--Vinogradov theorem of Chapter~4 up to level $N^{1/2-\varepsilon}$. A weighted sieve counts the $N-p$ with at most three prime factors, and the difficult step, the \emph{switching principle}, bounds those with exactly three by sifting a different list, of integers $N-p_1p_2p_3$, where the roles of the prime and the almost-prime are exchanged.

The second way adds information that is not of the form $|\mathcal A_d|$: bilinear sums, in which the sequence is tested against products $mn$ with both factors in ranges. These ``type II'' estimates break parity. With them, Friedlander and Iwaniec proved in 1998 that there are infinitely many primes of the form $a^2+b^4$, and Heath-Brown in 2001 that there are infinitely many of the form $a^3+2b^3$: thin sets, where the sieve alone has nothing to say. Vaughan's identity in Chapter~4 is the simplest machine of this kind.

The third way changes the question. Bounded gaps (Chapters~8 and~9) needs not one prime but \emph{two} primes among a few linear forms, and Maynard's sieve detects that some form is prime with positive probability without ever isolating a single one. Parity forbids finding primes, not finding them in tuples on average.

\begin{remark}[the twin constant]
Hardy and Littlewood conjectured $\pi_2(x)\sim2C_2\frac x{\log^2x}$, with $C_2=0.6601618\ldots$ Selberg's sieve applied to $\{n(n+2)\}$ gives $\pi_2(x)\le(16C_2+o(1))\frac x{\log^2x}$, eight times the conjecture (Exercise~\ref{ex:twins16}). Sifting instead $\{p+2:p\le x\}$, a sieve of dimension $1$ whose remainders are controlled by Bombieri--Vinogradov, Bombieri and Davenport obtained $8C_2$ in 1966. With much more machinery Wu reached $3.3996$ times the conjecture in 2004. By Corollary~\ref{cor:parity} and the think box above, no sieve can reach $1$ times it without new input on $\lambda$.
\end{remark}

\section{Artin's conjecture under GRH}

The decimal expansion of $\frac17=0.\overline{142857}$ has period $6=7-1$, the largest possible, because $10$ generates $(\Z/7\Z)^\times$: it is a \emph{primitive root} modulo $7$. Gauss tabulated such periods in the \emph{Disquisitiones}. In 1927, in a conversation with Helmut Hasse, Emil Artin conjectured that every integer $a$ other than $-1$ and the squares is a primitive root modulo infinitely many primes, and gave a heuristic density. For $a=2$ it is Artin's constant
\[
A=\prod_q\Bigl(1-\frac1{q(q-1)}\Bigr)=0.3739558\ldots,
\]
the product over all primes $q$. The heuristic is a sieve: $2$ fails to be a primitive root modulo $p$ exactly when, for some prime $q$, $p\equiv1\pmod q$ \emph{and} $2$ is a $q$-th power modulo $p$, two events of probabilities $\frac1{q-1}$ and $\frac1q$. In the late 1950s Derrick and Emma Lehmer's computations showed that for some $a$, such as $a=5$, the observed density disagreed with Artin's, and the correction (the events for different $q$ are not independent when $\Q(\sqrt a)$ sits inside a cyclotomic field) is part of the story told in Moree's survey. In 1967 Christopher Hooley proved the conjecture, with the corrected densities, assuming the Riemann hypothesis for the Dedekind zeta functions of certain number fields.

\begin{figure}[ht]
\centering
\fig{ch03-artin}
\caption{The proportion of primes $p\le x$ modulo which $a$ is a primitive root, for $a=2$ (charcoal), $3$ (rust), $5$ (ochre) and $10$ (coral), for $10^3\le x\le10^7$. At $x=10^7$ the proportions are $0.37391$, $0.37411$, $0.39394$ and $0.37449$. The lower dashed line is Artin's constant $A$, Hooley's density for $2$, $3$ and $10$; for $a=5\equiv1\pmod4$ the density is $\frac{20}{19}A=0.39364\ldots$, the upper dashed line, because $\sqrt5$ lies in $\Q(\zeta_5)$.}
\label{fig:artin}
\end{figure}

\subsection*{Kummer fields}

\begin{lemma}\label{lem:primroot}
Let $p$ be an odd prime. Then $2$ is a primitive root modulo $p$ if and only if there is no prime $q\mid p-1$ with $2^{(p-1)/q}\equiv1\pmod p$. For squarefree $k\mid p-1$, $2^{(p-1)/q}\equiv1$ for all primes $q\mid k$ if and only if $2^{(p-1)/k}\equiv1\pmod p$.
\end{lemma}
\begin{proof}
The order of $2$ divides $p-1$, and is a proper divisor exactly when it divides $\frac{p-1}q$ for some prime $q\mid p-1$. For the second claim, the order divides each $\frac{p-1}q$, $q\mid k$, if and only if it divides their greatest common divisor, which is $\frac{p-1}k$ as $k$ is squarefree.
\end{proof}

For squarefree $k$ let $\zeta_k=e^{2i\pi/k}$, $\alpha_k=2^{1/k}>0$, and $K_k=\Q(\zeta_k,\alpha_k)$. It is the splitting field of $X^k-2$, whose roots are the $\zeta_k^j\alpha_k$, so it is Galois over $\Q$.

\begin{lemma}\label{lem:kummerdeg}
For squarefree $k$, $[K_k:\Q]=k\varphi(k)$.
\end{lemma}
\begin{proof}
Let $F=\Q(\zeta_k)$, of degree $\varphi(k)$ with abelian Galois group (Volume~12, Chapter~5), and let $d=[K_k:F]\le k$. The minimal polynomial of $\alpha=\alpha_k$ over $F$ divides $\prod_j(X-\zeta_k^j\alpha)$, so its constant term is $\pm\zeta_k^b\alpha^d$ for some $b$, and $\alpha^d=2^{d/k}\in F$. With $g=\gcd(d,k)$ and $ud+vk=g$, also $2^{g/k}=(2^{d/k})^u2^v\in F$. If $g<k$, let $\ell$ be a prime dividing $\frac kg$; then $2^{1/\ell}=(2^{g/k})^{k/(g\ell)}\in F$. This is impossible:
\begin{itemize}
\item if $\ell$ is odd, $\Q(2^{1/\ell})$, a subfield of the abelian extension $F$, would be Galois over $\Q$, hence contain all the roots of the irreducible (Eisenstein at $2$) polynomial $X^\ell-2$, some of which are not real;
\item if $\ell=2$, then $\Q(\zeta_k)=\Q(\zeta_m)$ with $m$ the odd part of $k$, in which $2$ is unramified (Volume~12, Corollary~5.5), so $2$ would be unramified in its subfield $\Q(\sqrt2)$. The discriminant of $\Q(\sqrt2)$ is $8$, by the formula for quadratic fields in Volume~12, Chapter~1 (and the lattice drawn there). Theorem~2.15 of that volume is the ramification criterion, $p$ ramifies if and only if $p\mid d_K$, which then forbids $2$ from being unramified.
\end{itemize}
So $g=k$, $k\mid d$, and $d=k$.
\end{proof}

\begin{lemma}\label{lem:kummersplit}
Let $k$ be squarefree and $p$ a prime not dividing $2k$. Then $p$ is unramified in $K_k$, and $p$ splits completely in $K_k$ if and only if $p\equiv1\pmod k$ and $2^{(p-1)/k}\equiv1\pmod p$.
\end{lemma}
\begin{proof}
Let $\zeta=\zeta_k$, $\alpha=\alpha_k$, both algebraic integers, and $P$ a prime of $K_k$ above $p$. Every $\sigma\in\operatorname{Gal}(K_k/\Q)$ has $\sigma(\zeta)=\zeta^a$ and $\sigma(\alpha)=\zeta^b\alpha$ for some $a,b$, and is determined by them. Two facts: the $\zeta^j$, $0\le j<k$, are distinct modulo $P$, because $\prod_{j=1}^{k-1}(1-\zeta^j)=k$ (the value at $1$ of $\frac{X^k-1}{X-1}$) and $k\notin P$; and $\alpha\notin P$, because $\alpha^k=2\notin P$.
If $\sigma\in I_P$, then $\zeta^a\equiv\zeta$ forces $\zeta^a=\zeta$, and $(\zeta^b-1)\alpha\in P$ with $\alpha\notin P$ forces $\zeta^b=1$; so $\sigma=1$, $I_P$ is trivial and $e=1$.
Then $p$ splits completely if and only if $D_P$, generated by $\operatorname{Frob}_P$, is trivial. From $\operatorname{Frob}_P(\zeta)=\zeta^a\equiv\zeta^p$ we get $\zeta^a=\zeta^p$, so $\operatorname{Frob}_P$ fixes $\zeta$ if and only if $p\equiv1\pmod k$. In that case $\operatorname{Frob}_P(\alpha)=\zeta^b\alpha\equiv\alpha^p=\alpha\cdot2^{(p-1)/k}$, so $\zeta^b\equiv2^{(p-1)/k}\pmod P$, and $\operatorname{Frob}_P(\alpha)=\alpha$ if and only if $2^{(p-1)/k}\equiv1\pmod P$, that is modulo $P\cap\Z=p\Z$.
\end{proof}

\subsection*{Hooley's theorem}

For squarefree $k$ let $\pi_k(x)$ be the number of primes $p\le x$, $p\nmid2k$, that split completely in $K_k$. The only input from beyond this book is an effective Chebotarev theorem under the Riemann hypothesis.

\begin{theorem}[Lagarias--Odlyzko 1977, Serre 1981; quoted]\label{thm:effcheb}
Let $k$ be squarefree, and assume that all the zeros of the Dedekind zeta function $\zeta_{K_k}$ in the critical strip have real part $\frac12$. Then
\[
\pi_k(x)=\frac{\operatorname{li}(x)}{k\varphi(k)}+O\bigl(\sqrt x\log(kx)\bigr),
\]
with an absolute implied constant.
\end{theorem}

It is the prime number theorem for $K_k$ with the error term of the Riemann hypothesis, proved from the explicit formula for $\zeta_{K_k}$ as Volume~5 proved it for $\zeta$; the degree $k\varphi(k)$ comes from Lemma~\ref{lem:kummerdeg} and the logarithm of the discriminant of $K_k$ is $O(k\varphi(k)\log 2k)$. Hooley proved the estimate he needed directly in 1967.

\begin{theorem}[Hooley, 1967]\label{thm:hooley}
Assume the Riemann hypothesis for $\zeta_{K_k}$ for every squarefree $k$. Then the number $N(x)$ of primes $p\le x$ modulo which $2$ is a primitive root satisfies
\[
N(x)=A\operatorname{li}(x)+O\Bigl(\frac{x\log\log x}{\log^2x}\Bigr).
\]
In particular $2$ is a primitive root modulo infinitely many primes.
\end{theorem}
\begin{proof}
Call a prime $q$ \emph{bad} for $p$ if $q\mid p-1$ and $2^{(p-1)/q}\equiv1\pmod p$; by Lemma~\ref{lem:primroot}, $2$ is a primitive root modulo an odd $p$ exactly when no $q$ is bad for $p$. Let
\[
\xi_1=\tfrac16\log x,\qquad\xi_2=\frac{\sqrt x}{\log^2x},\qquad\xi_3=\sqrt x\log x,
\]
and $\mathcal N$ the number of primes $\xi_1<p\le x$ with no bad $q\le\xi_1$. Every bad $q$ is at most $p-1<x$, so
\[
\mathcal N-M_1-M_2-M_3\le N(x)\le\mathcal N+\xi_1,
\]
where $M_1$, $M_2$, $M_3$ count the $p\le x$ with a bad $q$ in $]\xi_1,\xi_2]$, $]\xi_2,\xi_3]$ and $]\xi_3,x[$.

\emph{Small $q$: the main term.} Let $Q=\prod_{q\le\xi_1}q$. By inclusion--exclusion over the set of bad $q\le\xi_1$,
\[
\mathcal N=\sum_{k\mid Q}\mu(k)\,\#\{\xi_1<p\le x:\text{every }q\mid k\text{ is bad for }p\}.
\]
By Lemmas~\ref{lem:primroot} and~\ref{lem:kummersplit}, and since $p>\xi_1$ does not divide $2k$, the condition means that $p$ splits completely in $K_k$; so the count is $\pi_k(x)+O(\xi_1)$. There are $2^{\pi(\xi_1)}\le2^{\xi_1}=x^{(\log2)/6}$ divisors $k\mid Q$, each at most $Q\le4^{\xi_1}\le x$. By Theorem~\ref{thm:effcheb},
\[
\mathcal N=\operatorname{li}(x)\sum_{k\mid Q}\frac{\mu(k)}{k\varphi(k)}+O\bigl(x^{(\log2)/6}\sqrt x\log x\bigr)=\operatorname{li}(x)\prod_{q\le\xi_1}\Bigl(1-\frac1{q(q-1)}\Bigr)+O(x^{0.62}).
\]
The product is $A$ divided by $\prod_{q>\xi_1}(1-\frac1{q(q-1)})\ge1-\sum_{n>\xi_1}\frac1{n(n-1)}\ge1-\frac1{\xi_1-1}$, so it is $A+O(\frac1{\log x})$, and $\mathcal N=A\operatorname{li}(x)+O(\frac x{\log^2x})$.

\emph{Medium $q$: GRH again.} If $q$ is bad for $p$, then $p$ splits completely in $K_q$. So
\[
M_1\le\sum_{\xi_1<q\le\xi_2}\pi_q(x)\le\operatorname{li}(x)\sum_{q>\xi_1}\frac1{q(q-1)}+O\bigl(\xi_2\sqrt x\log x\bigr)=O\Bigl(\frac x{\log^2x}\Bigr).
\]

\emph{Large $q$: the sieve.} Here we only use $q\mid p-1$. Since $\xi_3\le x^{3/4}$, Corollary~\ref{cor:btap} gives
\[
M_2\le\sum_{\xi_2<q\le\xi_3}\pi(x;q,1)\le\frac{2C_7x}{\log x}\sum_{\xi_2<q\le\xi_3}\frac1q.
\]
By Mertens, the last sum is $\log\frac{\log\xi_3}{\log\xi_2}+O(\frac1{\log x})$, and $\frac{\log\xi_3}{\log\xi_2}=\frac{\frac12\log x+\log\log x}{\frac12\log x-2\log\log x}=1+O(\frac{\log\log x}{\log x})$. So $M_2=O(\frac{x\log\log x}{\log^2x})$.

\emph{Very large $q$: an elementary count.} If $q>\xi_3$ is bad for $p$, then $m=\frac{p-1}q<\frac x{\xi_3}=\frac{\sqrt x}{\log x}$ and $2^m\equiv1\pmod p$: $p$ divides $2^m-1$ for some $m<\frac{\sqrt x}{\log x}$. As $2^m-1$ has fewer than $m$ prime factors, $M_3\le\sum_{m<\sqrt x/\log x}m\le\frac x{2\log^2x}$.

Adding up, $N(x)=A\operatorname{li}(x)+O(\frac{x\log\log x}{\log^2x})$.
\end{proof}

The proof is a sieve in which each range of $q$ needs its own tool: GRH for the $q$ up to $\frac{\sqrt x}{\log^2x}$, Brun--Titchmarsh for a short range near $\sqrt x$, and a triviality beyond. The gap the hypothesis fills is the range $\log x\ll q\ll\sqrt x$: without GRH, Chebotarev's theorem is effective only for fields of tiny degree.

\begin{remark}[without GRH]
For a general $a$, not $-1$ and not a square, Hooley's proof gives the density $A$ times a rational correction factor, which is $1$ unless the squarefree part of $a$ is $\equiv1\pmod4$ (for $a=5$ it is $\frac{20}{19}$; Exercise~\ref{ex:squares} explains the excluded $a$). Unconditionally, Gupta and Murty showed in 1984 that some explicit infinite set of $a$ satisfies the conjecture, and Heath-Brown deduced in 1986 that at most two primes $a$ fail it. So one of $2$, $3$, $5$ is a primitive root modulo infinitely many primes, but nobody knows which. Their method is a sieve (Chen's switching principle again), applied to the $p$ with $\frac{p-1}2$ having few prime factors. For $a=2$, as for every single given $a$, the conjecture remains open.
\end{remark}

\section{Exercises}

\begin{exercise}[$\star$]\label{ex:sg}
Let $q\equiv1\pmod4$ be a prime such that $p=2q+1$ is prime. Show that $2$ is a primitive root modulo $p$. What goes wrong for $q\equiv3\pmod4$? Explain why infinitely many such pairs $(q,p)$ would prove Artin's conjecture for $2$ unconditionally, and why sieve methods alone cannot provide them.
\end{exercise}

\begin{exercise}[$\star$]\label{ex:squares}
Show that a perfect square $b^2$ is a primitive root modulo no odd prime $p$, and that $-1$ is a primitive root modulo $3$ only. Where does Lemma~\ref{lem:kummerdeg} fail if $2$ is replaced by $b^2$, or by $-1$?
\end{exercise}

\begin{exercise}[$\star\star$]\label{ex:almostsg}
Show that there are $r$ and $c>0$ such that, for $x$ large, at least $c\frac x{\log^2x}$ integers $n\le x$ have $\Omega(n)\le r$ and $\Omega(2n+1)\le r$.
\end{exercise}

\begin{exercise}[$\star\star$]\label{ex:artinupper}
Without any hypothesis, show that the set of primes modulo which $2$ is a primitive root has upper Dirichlet density at most $A$: $\limsup_{\sigma\to1^+}\frac{\sum p^{-\sigma}}{\log\frac1{\sigma-1}}\le A$, the sum over these primes. \emph{(Use Volume~12, Theorem~8.23.)}
\end{exercise}

\begin{exercise}[$\star\star$]\label{ex:p2}
With the notation of Theorem~\ref{thm:parity}, show that $S(\mathcal A^-,\mathcal P,x^{1/3})=\pi(x)-\pi(x^{1/3})$ and that $S(\mathcal A^+,\mathcal P,x^{1/3})-1$ is the number of products $p_1p_2\le x$ of primes $x^{1/3}<p_1\le p_2$. Show that this number is $\asymp\frac x{\log x}$. What does this say about a sieve at level $x^{1-\varepsilon}$ applied to the integers up to $x$, with $z=x^{1/3}$?
\end{exercise}

\begin{exercise}[$\star\star\star$]\label{ex:twins16}
Apply Selberg's sieve (Volume~5, Theorem~8.2) to $\{n(n+2):n\le x\}$ as in Volume~5, Theorem~8.5, with $y=\frac{\sqrt x}{\log^3x}$, and show that
\[
G(y)=\sum_{d\le y}\mu^2(d)h(d)=\frac{\log^2y}{4C_2}+O(\log y),
\]
where $h(2)=1$ and $h(p)=\frac2{p-2}$ for odd $p$. Deduce $\pi_2(x)\le(16C_2+o(1))\frac x{\log^2x}$.
\end{exercise}

\section{Hints}

\begin{hint}{ex:sg}
$p\equiv3\pmod8$, and the prime factors of $p-1$ are $2$ and $q$. Use the second supplementary law for $\bigl(\frac2p\bigr)$.
\end{hint}

\begin{hint}{ex:squares}
Compute $(b^2)^{(p-1)/2}$. In the lemma, which step used that $2$ is not an $\ell$-th power of a real number in $F$?
\end{hint}

\begin{hint}{ex:almostsg}
Follow Corollary~\ref{cor:brunalmost}: count the solutions of $n(2n+1)\equiv0\pmod p$.
\end{hint}

\begin{hint}{ex:artinupper}
For fixed $\eta$, the primes with no bad $q\le\eta$ form a finite Boolean combination of the sets of primes splitting completely in the $K_k$, $k\mid\prod_{q\le\eta}q$.
\end{hint}

\begin{hint}{ex:p2}
An $n\le x$ whose prime factors all exceed $x^{1/3}$ has at most two of them. For the size, restrict to $x^{1/3}<p_1\le x^{0.45}$ and use Chebyshev's bounds and Mertens.
\end{hint}

\begin{hint}{ex:twins16}
Write $\mu^2h=\frac\tau n*b$, that is $\sum\mu^2(d)h(d)d^{-s}=\zeta(s+1)^2B(s)$, and show that $\sum|b(k)|k^{1/4}<\infty$. Then use $\sum_{m\le Y}\frac{\tau(m)}m=\frac12\log^2Y+O(\log Y+1)$.
\end{hint}

\section{Solutions}

\begin{solution}{ex:sg}
\idea Only two primes divide $p-1=2q$, so only two conditions need checking, and the quadratic character of $2$ settles the hard one.
\writeup Since $q\equiv1\pmod4$, $p=2q+1\equiv3\pmod8$, so $\bigl(\frac2p\bigr)=-1$ and $2^{(p-1)/2}\equiv-1\not\equiv1$. Also $2^{(p-1)/q}=4\not\equiv1\pmod p$, as $p\ge11$. By Lemma~\ref{lem:primroot}, $2$ is a primitive root. For $q\equiv3\pmod4$, $p\equiv7\pmod8$, $2$ is a square modulo $p$, and $2^{(p-1)/2}\equiv1$: $2$ is never a primitive root then.
Infinitely many such pairs would give infinitely many $p$ with $2$ as primitive root, without any hypothesis. But finding them means finding primes $q$ for which $2q+1$ is also prime, a problem of the same kind as twin primes; by Corollary~\ref{cor:parity}, a sieve can give an upper bound of the right order (Exercise~\ref{ex:almostsg} gives the almost-prime version from below) but never a single pair.
\end{solution}

\begin{solution}{ex:squares}
\idea A square is a square: it lies in the subgroup of index $2$, which a generator never does.
\writeup If $p\mid b$, then $b^2\equiv0$ is not even invertible. For an odd prime $p\nmid b$, $(b^2)^{(p-1)/2}=b^{p-1}\equiv1$, so the order of $b^2$ divides $\frac{p-1}2<p-1$. So $b^2$ is never a primitive root modulo an odd prime. The order of $-1$ is $2$, so it is a primitive root only when $p-1=2$, that is $p=3$.
In Lemma~\ref{lem:kummerdeg} with $b^2$ in place of $2$, the step $\ell=2$ fails: $\sqrt{b^2}=b\in F$, so $[K_2:\Q]=1$, not $2$, and every $p$ ``splits completely'' in $K_2$: the prime $q=2$ is bad for every $p$, and the density is $0$. With $-1$, the step for odd $\ell$ fails: $-1=(-1)^\ell$ is an $\ell$-th power in $\Q$, so $K_\ell=\Q(\zeta_\ell)$ has degree $\ell-1$, not $\ell(\ell-1)$.
\end{solution}

\begin{solution}{ex:almostsg}
\idea The polynomial $n(2n+1)$ has two roots modulo every odd prime and one modulo $2$, exactly like $n(n+2)$; the sieve sees no difference.
\writeup For an odd prime $p$, $n(2n+1)\equiv0\pmod p$ means $n\equiv0$ or $n\equiv-\frac12$, two classes; modulo $2$, $2n+1$ is odd, so only $n\equiv0$. So $\rho$ is as in Corollary~\ref{cor:brunalmost}: $|\mathcal A_d|=\frac{\rho(d)}dx+r_d$, $|r_d|\le\tau(d)$, and $(\Omega)$ holds with $\kappa=2$, $K=3$. The proof of Corollary~\ref{cor:brunalmost} goes through word for word, giving $\gg\frac x{\log^2x}$ integers $n\le x$ such that $n$ and $2n+1$ have all their prime factors above $z=x^{1/(2s_1)}$; as $2n+1\le3x$, $\Omega(n)$ and $\Omega(2n+1)$ are at most $\lfloor2s_1\rfloor+2$.
\end{solution}

\begin{solution}{ex:artinupper}
\idea Throwing away the large $q$ only enlarges the set, and for each finite set of $q$ Chebotarev's theorem gives the exact density of the larger set.
\writeup Fix $\eta\ge2$, let $Q_\eta=\prod_{q\le\eta}q$, and let $S_\eta$ be the set of primes $p>\eta$ with no bad $q\le\eta$. The set $S$ of primes modulo which $2$ is a primitive root satisfies $S\setminus[1,\eta]\subseteq S_\eta$. For $k\mid Q_\eta$ and $p>\eta$, ``every $q\mid k$ is bad'' means that $p$ splits completely in $K_k$ (Lemmas~\ref{lem:primroot} and~\ref{lem:kummersplit}), and these primes have Dirichlet density $\frac1{[K_k:\Q]}=\frac1{k\varphi(k)}$, by Volume~12, Theorem~8.23 applied to the class of the identity, and Lemma~\ref{lem:kummerdeg}. Inclusion--exclusion is a finite sum of indicator functions, so
\[
\sum_{p\in S_\eta}p^{-\sigma}=\sum_{k\mid Q_\eta}\mu(k)\sum_{\substack{p>\eta\\p\text{ splits in }K_k}}p^{-\sigma}=\Bigl(\sum_{k\mid Q_\eta}\frac{\mu(k)}{k\varphi(k)}+o(1)\Bigr)\log\frac1{\sigma-1}
\]
as $\sigma\to1^+$. The sum is $\prod_{q\le\eta}(1-\frac1{q(q-1)})$. So the upper Dirichlet density of $S$ is at most this product for every $\eta$, hence at most its limit $A$.
\end{solution}

\begin{solution}{ex:p2}
\idea At $z=x^{1/3}$ the survivors are primes or products of two primes, and $\lambda$ sorts them perfectly into the two lists.
\writeup If $n\le x$ has all its prime factors above $x^{1/3}$, it has at most two. So the survivors in $\mathcal A^-$ ($\Omega$ odd) are the primes in $]x^{1/3},x]$, and those in $\mathcal A^+$ ($\Omega$ even) are $1$ and the $p_1p_2\le x$ with $x^{1/3}<p_1\le p_2$.
For each prime $p_1\in\left]x^{1/3},x^{0.45}\right]$, the number of $p_2$ in $[p_1,\frac x{p_1}]$ is $\pi(\frac x{p_1})-\pi(p_1)+1\gg\frac x{p_1\log x}$, by Chebyshev's bounds, since $\frac x{p_1}\ge x^{0.55}$ is much larger than $p_1$. Summing, the count is $\gg\frac x{\log x}\sum_{x^{1/3}<p\le x^{0.45}}\frac1p\gg\frac x{\log x}$ by Mertens ($\log\frac{0.45}{1/3}>0$). It is at most $\sum_{p_1\le\sqrt x}\pi(\frac x{p_1})\ll\frac x{\log x}\sum_{x^{1/3}<p\le\sqrt x}\frac1p\ll\frac x{\log x}$. (With the prime number theorem one finds $\sim\log2\cdot\frac x{\log x}$.)
So at $z=x^{1/3}$ the integers up to $x$ have a positive proportion of survivors of each kind, and the lists $\mathcal A^\pm$, which have the same sieve data, keep one kind each. A sieve using only the data $|\mathcal A_d|$, $d<x^{1-\varepsilon}$, cannot tell whether its survivors are primes or products of two primes. It can, however, count survivors with $\Omega\le2$ from below, which is the opening Chen used.
\end{solution}

\begin{solution}{ex:twins16}
\idea $h(p)$ is $\frac2p$ up to $O(p^{-2})$, so $\mu^2h$ behaves like $\frac{\tau(n)}n$ times a convergent Euler product, whose value at the point that matters is $\frac1{2C_2}$.
\writeup Let $f=\mu^2h$, with $F(s)=\sum f(n)n^{-s}=\prod_p(1+h(p)p^{-s})$, and write $F(s)=\zeta(s+1)^2B(s)$ with $B(s)=\prod_p(1+h(p)p^{-s})(1-p^{-1-s})^2=\sum b(k)k^{-s}$. The coefficients $b$ are multiplicative, $b(p^j)=0$ for $j\ge4$, and
\[
b(p)=h(p)-\frac2p,\quad b(p^2)=\frac1{p^2}-\frac{2h(p)}p,\quad b(p^3)=\frac{h(p)}{p^2}.
\]
For odd $p$, $b(p)=\frac4{p(p-2)}$, so $|b(p^j)|\ll p^{-2}$ for $j=1,2,3$; for $p=2$, $b(2)=0$, $b(4)=-\frac34$, $b(8)=\frac14$. Hence $\sum_j|b(p^j)|p^{j/4}\ll p^{-3/2}$, the Euler product converges, and $\sum_k|b(k)|k^{1/4}<\infty$. In particular $\sum_k|b(k)|(1+\log k)^2<\infty$ and $\sum_{k>y}|b(k)|\le y^{-1/4}\sum_k|b(k)|k^{1/4}=O(y^{-1/4})$.
Comparing coefficients in $F(s)=B(s)\sum_m\tau(m)m^{-1-s}$ gives $f(n)=\sum_{km=n}b(k)\frac{\tau(m)}m$. So
\[
G(y)=\sum_{k\le y}b(k)T\Bigl(\frac yk\Bigr),\qquad T(Y)=\sum_{m\le Y}\frac{\tau(m)}m.
\]
Now $T(Y)=\sum_{a\le Y}\frac1a\sum_{c\le Y/a}\frac1c=\sum_{a\le Y}\frac1a\bigl(\log\frac Ya+O(1)\bigr)=\log Y\sum_{a\le Y}\frac1a-\sum_{a\le Y}\frac{\log a}a+O(\log Y+1)$, and with $\sum_{a\le Y}\frac1a=\log Y+O(1)$, $\sum_{a\le Y}\frac{\log a}a=\frac12\log^2Y+O(1)$ (compare with the integral), $T(Y)=\frac12\log^2Y+O(\log Y+1)$. Therefore
\[
G(y)=\frac12\sum_{k\le y}b(k)(\log y-\log k)^2+O\Bigl(\sum_k|b(k)|(\log y+1)\Bigr)=\frac{B(0)}2\log^2y+O(\log y),
\]
where expanding the square, the terms with $\log k$ are $O(\log y)$, and replacing $\sum_{k\le y}b(k)$ by $B(0)=\sum_kb(k)$ costs $O(y^{-1/4}\log^2y)$. Finally
\[
B(0)=\prod_p(1+h(p))\Bigl(1-\frac1p\Bigr)^2=\frac12\prod_{p>2}\frac{(p-1)^2}{p(p-2)}=\frac1{2C_2},
\]
since $(1+h(2))\frac14=\frac12$ and, for odd $p$, $(1+\frac2{p-2})(1-\frac1p)^2=\frac{(p-1)^2}{p(p-2)}=\bigl(1-\frac1{(p-1)^2}\bigr)^{-1}$. So $G(y)=\frac{\log^2y}{4C_2}+O(\log y)$.
With $z=y$, every squarefree $d\le y$ divides $P(z)$, and as in Volume~5, Theorem~8.5, the remainder is at most $\bigl(\sum_{d\le y}\tau(d)\bigr)^2\le y^2(1+\log y)^2\le\frac x{\log^3x}$ for large $x$. Since $\log y=\frac12\log x-3\log\log x$,
\[
\pi_2(x)\le z+\frac x{G(y)}+\frac x{\log^3x}=\frac{4C_2x}{\log^2y}(1+o(1))=(16C_2+o(1))\frac x{\log^2x}.
\]
\end{solution}

\section*{Sources and further reading}

\begin{itemize}
\item \emph{Brun's sieve.} V.~Brun, ``La s\'erie $\frac15+\frac17+\frac1{11}+\frac1{13}+\cdots$ o\`u les d\'enominateurs sont nombres premiers jumeaux est convergente ou finie'', \emph{Bull.\ Sci.\ Math.}\ 43 (1919), 100--104 and 124--128; ``Le crible d'Eratosthène et le théorème de Goldbach'', \emph{Videnskapsselskapets Skrifter, Kristiania} I, no.~3 (1920). The fundamental lemma in this form, with Brun's ranges and Bonferroni, follows J.~Friedlander and H.~Iwaniec, \emph{Opera de Cribro}, AMS Colloquium Publications 57, 2010, Chapter~6, and H.~Halberstam and H.-E.~Richert, \emph{Sieve Methods}, Academic Press, 1974, Chapter~2.
\item \emph{Parity.} A.~Selberg, ``Lectures on sieves'', in \emph{Collected Papers}, Vol.~II, Springer, 1991, 65--247. E.~Bombieri, ``The asymptotic sieve'', \emph{Rend.\ Accad.\ Naz.\ XL} (5) 1/2 (1975/76), 243--269. T.~Tao, ``The logarithmically averaged Chowla and Elliott conjectures for two-point correlations'', \emph{Forum Math.\ Pi} 4 (2016), e8.
\item \emph{Beyond parity.} J.-R.~Chen, ``On the representation of a larger even integer as the sum of a prime and the product of at most two primes'', \emph{Sci.\ Sinica} 16 (1973), 157--176. J.~Friedlander and H.~Iwaniec, ``The polynomial $X^2+Y^4$ captures its primes'', \emph{Ann.\ of Math.}\ 148 (1998), 945--1040. D.~R.~Heath-Brown, ``Primes represented by $x^3+2y^3$'', \emph{Acta Math.}\ 186 (2001), 1--84.
\item \emph{Twin prime constants.} E.~Bombieri and H.~Davenport, ``Small differences between prime numbers'', \emph{Proc.\ Roy.\ Soc.\ Ser.\ A} 293 (1966), 1--18. J.~Wu, ``Chen's double sieve, Goldbach's conjecture and the twin prime problem'', \emph{Acta Arith.}\ 114 (2004), 215--273.
\item \emph{Artin's conjecture.} C.~Hooley, ``On Artin's conjecture'', \emph{J.\ Reine Angew.\ Math.}\ 225 (1967), 209--220. J.~C.~Lagarias and A.~M.~Odlyzko, ``Effective versions of the Chebotarev density theorem'', in \emph{Algebraic Number Fields}, A.~Fröhlich (ed.), Academic Press, 1977, 409--464. J.-P.~Serre, ``Quelques applications du th\'eor\`eme de densit\'e de Chebotarev'', \emph{Publ.\ Math.\ IH\'ES} 54 (1981), 123--201. R.~Gupta and M.~R.~Murty, ``A remark on Artin's conjecture'', \emph{Invent.\ Math.}\ 78 (1984), 127--130. D.~R.~Heath-Brown, ``Artin's conjecture for primitive roots'', \emph{Quart.\ J.\ Math.\ Oxford} (2) 37 (1986), 27--38. P.~Moree, ``Artin's primitive root conjecture: a survey'', \emph{Integers} 12A (2012), A13: the history, including Artin's conversation with Hasse and the Lehmers' computations.
\end{itemize}
